\chapter{数学预备}
\label{chap:advanced-mathematical-foundations}

物理公式中的对象各有明确类型。波函数是向量，哈密顿量是线性算符，对称操作组成群，参数依赖的本征空间组成向量丛。若不先分清这些对象，后面的计算即使符号无误，也容易把“选了一组基”误认为“改变了物理状态”，或把局部坐标公式误认为全局结论。这里从定义开始建立全书反复使用的数学语言。

\section{集合、映射与商集}

写 \(x\in X\) 表示 \(x\) 是集合 \(X\) 的元素。映射
\[
f:X\longrightarrow Y
\]
给每个 \(x\in X\) 指定唯一的 \(f(x)\in Y\)。若不同的 \(x\) 总有不同的像，\(f\) 称为单射；若每个 \(y\in Y\) 都等于某个 \(f(x)\)，\(f\) 称为满射；既单又满的映射称为双射。双射存在逆映射。

关系 \(\sim\) 若满足
\[
x\sim x,\qquad
x\sim y\Rightarrow y\sim x,\qquad
x\sim y,\ y\sim z\Rightarrow x\sim z,
\]
便称为等价关系。元素 \(x\) 的等价类为
\[
[x]=\{y\in X:y\sim x\},
\]
全体等价类组成商集 \(X/{\sim}\)。例如角度 \(\theta\) 与
\(\theta+2\pi n\) 表示同一方向，所以圆周可写成
\(\mathbb R/(2\pi\mathbb Z)\)。

\begin{proposition}[等价类要么相同，要么不相交]
若 \(\sim\) 是等价关系，则任意两个等价类 \([x]\) 与 \([y]\) 或者相等，或者没有公共元素。
\end{proposition}

\begin{proof}
若两类有公共元素 \(z\)，则 \(z\sim x\) 且 \(z\sim y\)，从而
\(x\sim y\)。对任意 \(u\in[x]\)，有 \(u\sim x\sim y\)，故
\(u\in[y]\)，所以 \([x]\subseteq[y]\)。交换 \(x,y\) 得到反向包含，于是 \([x]=[y]\)。
\end{proof}

\section{向量空间、线性映射与商空间}

\begin{definitionplain}[向量空间]
在数域 \(\mathbb F\) 上，集合 \(V\) 若定义了向量加法和数乘，并满足加法结合律、交换律、零元与负元存在，以及数乘的结合律和分配律，就称为 \(\mathbb F\) 上的向量空间。本书主要取
\(\mathbb F=\mathbb R\) 或 \(\mathbb C\)。
\end{definitionplain}

一组向量 \(e_1,\ldots,e_n\) 若线性无关，并且每个 \(v\in V\) 都能唯一写成 \(v=\sum_i v^ie_i\)，便是一组基。数 \(v^i\) 是 \(v\) 在这组基下的坐标。换基只改变坐标，不改变向量。

映射 \(T:V\to W\) 若满足
\[
T(av+bw)=aT(v)+bT(w),
\]
便是线性映射。它的核与像分别为
\[
\ker T=\{v:T(v)=0\},\qquad
\operatorname{im}T=\{T(v):v\in V\}.
\]

\begin{theorem}[秩--零化度定理]
设 \(V\) 有限维，\(T:V\to W\) 线性，则
\[
\dim V=\dim\ker T+\dim\operatorname{im}T.
\]
\end{theorem}

\begin{proof}
取 \(\ker T\) 的基 \(e_1,\ldots,e_k\)，再扩充为 \(V\) 的基
\(e_1,\ldots,e_n\)。向量
\(T(e_{k+1}),\ldots,T(e_n)\) 张成 \(\operatorname{im}T\)。若它们有线性关系，则相应的
\(\sum_{i=k+1}^na_ie_i\) 属于 \(\ker T\)，与整组基线性无关矛盾，所以它们也是线性无关的。因此
\(\dim\operatorname{im}T=n-k\)，结论成立。
\end{proof}

若 \(U\subseteq V\) 是子空间，可规定 \(v\sim w\) 当且仅当
\(v-w\in U\)。所得商空间 \(V/U\) 的元素是陪集
\([v]=v+U\)。

\begin{proposition}[商空间维数]
有限维情形下，
\[
\dim(V/U)=\dim V-\dim U.
\]
\end{proposition}

\begin{proof}
取 \(U\) 的基 \(e_1,\ldots,e_k\)，扩充为 \(V\) 的基
\(e_1,\ldots,e_n\)。在商空间中，前 \(k\) 个基向量的等价类为零，而
\([e_{k+1}],\ldots,[e_n]\) 线性无关并张成 \(V/U\)，所以维数为
\(n-k\)。
\end{proof}

\section{内积、伴随与谱分解}

复向量空间上的内积 \(\langle u,v\rangle\) 对第二个变量线性，对第一个变量共轭线性，并满足
\[
\langle u,v\rangle=\langle v,u\rangle^*,\qquad
\langle v,v\rangle\ge0.
\]
等号只在 \(v=0\) 时成立。算符 \(A\) 的伴随 \(A^\dagger\) 由
\[
\langle u,Av\rangle=\langle A^\dagger u,v\rangle
\]
定义。若 \(A=A^\dagger\)，则 \(A\) 自伴；若
\(A^\dagger A=I\)，则 \(A\) 酉。

\begin{proposition}[自伴算符的基本谱性质]
自伴算符的本征值为实数，属于不同本征值的本征向量彼此正交。
\end{proposition}

\begin{proof}
若 \(Av=\lambda v\)，则
\[
\lambda\langle v,v\rangle
=\langle v,Av\rangle
=\langle Av,v\rangle
=\lambda^*\langle v,v\rangle,
\]
所以 \(\lambda=\lambda^*\)。若 \(Av=\lambda v\)、\(Aw=\mu w\)，则
\[
\mu\langle v,w\rangle
=\langle v,Aw\rangle
=\langle Av,w\rangle
=\lambda\langle v,w\rangle.
\]
当 \(\lambda\ne\mu\) 时只能有 \(\langle v,w\rangle=0\)。
\end{proof}

有限维谱定理说明，每个复自伴矩阵都能由酉变换对角化。证明可先从特征多项式取一个本征向量，再证明它的正交补在算符作用下保持不变，最后对维数归纳。

\section{群、表示与 Schur 引理}

\begin{definitionplain}[群]
集合 \(G\) 配有乘法。若乘法满足结合律，存在单位元 \(e\)，且每个
\(g\) 都有逆元 \(g^{-1}\)，则 \(G\) 称为群。若还有
\(gh=hg\)，则群是交换群。
\end{definitionplain}

\begin{definitionplain}[表示]
群 \(G\) 在向量空间 \(V\) 上的表示，是群同态
\[
D:G\longrightarrow GL(V),
\qquad D(gh)=D(g)D(h).
\]
若每个 \(D(g)\) 都酉，则称为酉表示。
\end{definitionplain}

子空间 \(W\subseteq V\) 若满足 \(D(g)W\subseteq W\) 对所有
\(g\in G\) 成立，称为不变子空间。除 \(0,V\) 外没有不变子空间的表示称为不可约表示。线性映射 \(A:V\to W\) 若满足
\[
AD_V(g)=D_W(g)A,
\]
称为交织算符。

\begin{theorem}[Schur 引理]
设 \(V,W\) 是有限维复不可约表示。交织算符 \(A:V\to W\) 或为零，或为同构。若 \(V=W\)，则 \(A=cI\)。
\end{theorem}

\begin{proof}
\(\ker A\) 是 \(V\) 的不变子空间，\(\operatorname{im}A\) 是 \(W\) 的不变子空间。不可约性说明非零 \(A\) 的核只能为零，像只能为
\(W\)，所以它是同构。若 \(V=W\)，复矩阵 \(A\) 有某个本征值
\(c\)。交织算符 \(A-cI\) 有非零核，因此只能等于零。
\end{proof}

表示的特征标是 \(\chi(g)=\operatorname{Tr}D(g)\)。相似矩阵迹相同，而
\(D(hgh^{-1})=D(h)D(g)D(h)^{-1}\)，所以特征标在同一共轭类上相等。

\section{Lie 群、生成元与指数映射}

若群 \(G\) 同时是光滑流形，且乘法与取逆光滑，\(G\) 称为 Lie 群。单位元处的切空间
\(\mathfrak g=T_eG\) 称为 Lie 代数。

对矩阵 Lie 群，经过单位元的一参数子群满足
\(g(t+s)=g(t)g(s)\)。令 \(X=g'(0)\)，对 \(s\) 在零点求导得
\[
g'(t)=g(t)X,\qquad g(0)=I.
\]
这个初值问题的唯一解为
\[
g(t)=\exp(tX)=\sum_{n=0}^{\infty}\frac{t^nX^n}{n!}.
\]
生成元的交换子 \([X,Y]=XY-YX\) 测量两个无穷小变换交换次序后的最低阶差异。

\begin{proposition}[正交群的 Lie 代数]
\[
SO(d)=\{R:R^{\mathsf T}R=I,\ \det R=1\}
\]
的 Lie 代数是全体实反对称矩阵：
\[
\mathfrak{so}(d)=\{X:X^{\mathsf T}=-X\}.
\]
\end{proposition}

\begin{proof}
取 \(R(t)\in SO(d)\)、\(R(0)=I\)、\(R'(0)=X\)。对
\(R(t)^{\mathsf T}R(t)=I\) 求导，得到
\(X^{\mathsf T}+X=0\)。反过来，若 \(X^{\mathsf T}=-X\)，则
\[
(\ee^{tX})^{\mathsf T}\ee^{tX}
=\ee^{-tX}\ee^{tX}=I,
\]
并且行列式连续且在 \(t=0\) 等于 \(1\)，所以
\(\ee^{tX}\in SO(d)\)。
\end{proof}

\section{流形、切向量与微分形式}

曲线、曲面和转动群的共同点，不是它们都能嵌入某个更大的欧氏空间，而是每个点附近都可以用若干实数作坐标。球面需要至少两张坐标图；用经纬度覆盖时，极点处的经度失去意义，但球面本身并没有在那里产生奇点。微分几何首先要把这种“坐标失效”和真正的几何奇点区分开。

\begin{definitionplain}[光滑流形]
一个 \(n\) 维光滑流形 \(M\) 是 Hausdorff、第二可数的拓扑空间，并带有一族坐标图
\[
(U_\alpha,x_\alpha),\qquad
x_\alpha:U_\alpha\longrightarrow x_\alpha(U_\alpha)\subseteq\mathbb R^n.
\]
各 \(U_\alpha\) 覆盖 \(M\)，每个 \(x_\alpha\) 都是到 \(\mathbb R^n\) 中开集的同胚；若 \(U_\alpha\cap U_\beta\ne\varnothing\)，则坐标变换
\[
x_\beta\circ x_\alpha^{-1}:
x_\alpha(U_\alpha\cap U_\beta)\longrightarrow
x_\beta(U_\alpha\cap U_\beta)
\]
光滑，并有光滑逆映射。满足这些条件的一整族相容坐标图称为光滑图册。
\end{definitionplain}

Hausdorff 条件保证两个不同点能由不相交邻域分开，第二可数条件排除过大的病态空间。物理中常见的开集、球面、环面、矩阵群和由正则约束定义的曲面都满足这两条条件。真正承担微分运算的是坐标变换的光滑性：同一个函数在两张图中写法不同，但链式法则保证导数能够相互转换。

单位圆 \(S^1=\{(x,y)\in\mathbb R^2:x^2+y^2=1\}\) 已经说明为什么一张全局坐标图通常不存在。删去一点以后可用角度区间参数化；另删一点可得到第二张图。重叠处的两个角坐标相差 \(2\pi\) 的整数倍，转移函数光滑。圆周也可以写成商空间 \(\mathbb R/(2\pi\mathbb Z)\)，这与前面等价关系的定义完全一致。

流形上的光滑函数由坐标来检验。若 \(f:M\to\mathbb R\)，则要求每个局部表达
\(f\circ x_\alpha^{-1}\) 都是普通的光滑函数。若 \(F:M\to N\) 是两个流形间的映射，则在任意相交坐标图中要求
\[
y_\beta\circ F\circ x_\alpha^{-1}
\]
光滑。这个定义不依赖图册中的某一张图，因为不同图之间的转换本身已经光滑。

在欧氏空间里，点 \(p\) 处的方向是一个箭头；在抽象流形上，没有预先给定的环境空间可供放置箭头。可以改用过 \(p\) 的曲线来定义方向。设 \(\gamma:(-\epsilon,\epsilon)\to M\) 且 \(\gamma(0)=p\)。若对每个光滑函数 \(f\) 都有
\[
\left.\frac{\dd}{\dd t}f(\gamma_1(t))\right|_{t=0}
=
\left.\frac{\dd}{\dd t}f(\gamma_2(t))\right|_{t=0},
\]
便把 \(\gamma_1\) 与 \(\gamma_2\) 看作同一个切向量。于是切向量 \(X\) 可以直接视为作用在光滑函数上的导子，
\[
X[f]=\left.\frac{\dd}{\dd t}f(\gamma(t))\right|_{t=0},
\qquad
X[fg]=X[f]g(p)+f(p)X[g].
\]
所有切向量组成 \(n\) 维向量空间 \(T_pM\)，称为 \(p\) 点的切空间。

取局部坐标 \((x^1,\ldots,x^n)\)。普通多元微积分的链式法则给出
\[
X[f]=X^\mu\frac{\partial f}{\partial x^\mu}(p),
\qquad
X=X^\mu\left.\frac{\partial}{\partial x^\mu}\right|_p.
\]
因此 \(\partial/\partial x^\mu|_p\) 是 \(T_pM\) 的一组基。更换坐标 \(y^a=y^a(x)\) 时，
\begin{equation}
\frac{\partial}{\partial x^\mu}
=\frac{\partial y^a}{\partial x^\mu}
\frac{\partial}{\partial y^a},
\qquad
X'^a=\frac{\partial y^a}{\partial x^\mu}X^\mu.
\label{eq:tangent-coordinate-transform-primer}
\end{equation}
坐标分量会变，导子 \(X\) 本身不变。

光滑映射 \(F:M\to N\) 在每个点诱导线性映射
\[
F_{*p}:T_pM\longrightarrow T_{F(p)}N,
\qquad
(F_{*p}X)[h]=X[h\circ F],
\]
称为微分或推出。若在坐标中写 \(F^a= y^a\circ F\)，则
\[
(F_{*p}X)^a=\frac{\partial F^a}{\partial x^\mu}(p)X^\mu,
\]
这正是 Jacobian 矩阵作用在方向导数上。复合映射满足
\((G\circ F)_{*p}=G_{*F(p)}\circ F_{*p}\)，所以推出是链式法则的无坐标写法。

切空间的对偶空间 \(T_p^*M\) 称为余切空间。坐标函数的微分 \(\dd x^\mu\) 由
\[
\dd x^\mu\!\left(\frac{\partial}{\partial x^\nu}\right)
=\delta^\mu_{\nu}
\]
定义，因此 \(\{\dd x^\mu\}\) 是余切基。对光滑函数 \(f\)，其微分 \(\dd f\in T_p^*M\) 满足
\[
(\dd f)_p(X)=X[f],
\qquad
\dd f=\frac{\partial f}{\partial x^\mu}\dd x^\mu.
\]
由式 \eqnrefs{eq:tangent-coordinate-transform-primer} 可知余切基按逆 Jacobian 变换；这正是上指标和下指标变换规律不同的来源。

在每一点平滑选取一个切向量得到向量场。两个向量场的 Lie 括号定义为
\[
[X,Y][f]=X[Y[f]]-Y[X[f]].
\]
在局部坐标中，若 \(X=X^\mu\partial_\mu\)、\(Y=Y^\mu\partial_\mu\)，则
\[
[X,Y]
=\left(X^\mu\partial_\mu Y^\nu-Y^\mu\partial_\mu X^\nu\right)\partial_\nu.
\]
坐标基满足 \([\partial_\mu,\partial_\nu]=0\)，一般移动标架则未必对易。

\begin{definitionplain}[微分形式]
点 \(p\) 处的 \(k\)-形式是 \(T_pM\) 上的反对称 \(k\) 重线性函数。把它随 \(p\) 光滑选取，便得到流形上的微分 \(k\)-形式。局部可写成
\[
\omega=\frac1{k!}\omega_{\mu_1\cdots\mu_k}
\dd x^{\mu_1}\wedge\cdots\wedge\dd x^{\mu_k},
\]
其中系数对全部指标完全反对称。
\end{definitionplain}

楔积把一个 \(k\)-形式和一个 \(\ell\)-形式合成 \((k+\ell)\)-形式，并满足
\[
\alpha\wedge\beta=(-1)^{k\ell}\beta\wedge\alpha.
\]
特别地，\(\dd x^\mu\wedge\dd x^\mu=0\)。外微分先在函数上由通常的全微分定义，再由
\[
\dd^2=0,
\qquad
\dd(\alpha\wedge\beta)
=\dd\alpha\wedge\beta+(-1)^k\alpha\wedge\dd\beta
\]
唯一延拓到所有形式。局部表达为
\[
\dd\omega
=\frac1{k!}\partial_\nu\omega_{\mu_1\cdots\mu_k}
\dd x^\nu\wedge\dd x^{\mu_1}\wedge\cdots\wedge\dd x^{\mu_k}.
\]
式 \(\dd^2=0\) 来自二阶偏导的对称性与楔积的反对称性。

映射 \(F:M\to N\) 还能把 \(N\) 上的微分形式拉回 \(M\)。对 \(k\)-形式 \(\omega\)，定义
\[
(F^*\omega)_p(X_1,\ldots,X_k)
=\omega_{F(p)}(F_{*p}X_1,\ldots,F_{*p}X_k).
\]
拉回满足 \(F^*(\alpha\wedge\beta)=F^*\alpha\wedge F^*\beta\) 和
\(F^*(\dd\omega)=\dd(F^*\omega)\)。积分变量替换、曲面上的面积元以及电磁场强在坐标变换下的写法，都使用这个运算。

在可定向流形上选择一致的方向以后，紧支撑的最高次形式可以积分。若 \(M\) 是带边界的定向 \(n\) 维流形，\(\omega\) 是紧支撑的 \((n-1)\)-形式，则 Stokes 定理为
\begin{equation}
\int_M\dd\omega=\int_{\partial M}\omega.
\label{eq:general-stokes-primer}
\end{equation}
边界方向取“外法向优先”的约定。把流形分割成坐标小块后，相邻小块在公共边界上诱导相反方向，内部积分两两抵消；每个小块内的结论则退化为多元微积分基本定理。微积分基本定理、Green 公式、Gauss 公式和经典 Stokes 公式因而都是式 \eqnrefs{eq:general-stokes-primer} 的不同维数版本。

\section{向量丛、联络与曲率}

流形把每一点附近的几何与 \(\mathbb R^n\) 比较；向量丛则在每一点上再附着一个向量空间。最简单的情形是直积 \(M\times\mathbb C^r\)，但电子本征态随核坐标变化时，通常只能在局部连续地选取一组本征矢，不能预先假定存在全局统一的基。

\begin{definitionplain}[向量丛与截面]
秩为 \(r\) 的复向量丛由流形 \(E,M\) 和光滑满射 \(\pi:E\to M\) 构成。每个纤维 \(E_p=\pi^{-1}(p)\) 是 \(r\) 维复向量空间；对每个 \(p\in M\)，存在邻域 \(U\) 和微分同胚
\[
\Phi_U:\pi^{-1}(U)\longrightarrow U\times\mathbb C^r,
\]
使 \(\Phi_U\) 在每个纤维上都是线性同构，并且投影到第一个因子后等于 \(\pi\)。满足 \(\pi(s(p))=p\) 的光滑映射 \(s:M\to E\) 称为截面。
\end{definitionplain}

切丛 \(TM=\bigsqcup_{p\in M}T_pM\) 是实秩 \(n\) 的向量丛，余切丛 \(T^*M\) 是其对偶丛。量子力学中若哈密顿量 \(H(R)\) 的某个隔离能级在每个核构型 \(R\) 上有 \(r\) 维本征空间，这些本征空间在能隙不闭合的区域内组成秩 \(r\) 的本征态丛。

在 \(U_\alpha\) 上逐点选取纤维基
\(e_{\alpha1},\ldots,e_{\alpha r}\)，称为局部标架。把它们排成行向量 \(e_\alpha\)，则重叠区上存在光滑可逆矩阵
\[
e_\beta=e_\alpha U_{\alpha\beta},
\qquad U_{\alpha\beta}:U_\alpha\cap U_\beta\to GL(r,\mathbb C).
\]
同一截面写成
\(s=e_\alpha\psi_\alpha=e_\beta\psi_\beta\)，故
\[
\psi_\beta=U_{\alpha\beta}^{-1}\psi_\alpha.
\]
三个坐标邻域相交时，同一个标架变换可以分两步完成，所以转移函数必须满足
\[
U_{\alpha\alpha}=I,
\qquad U_{\alpha\beta}U_{\beta\gamma}=U_{\alpha\gamma}.
\]
这就是余圈条件。反过来，一族满足余圈条件的转移函数可以把各个 \(U_\alpha\times\mathbb C^r\) 粘合成向量丛。

不同纤维 \(E_p\) 与 \(E_q\) 是不同的向量空间，所以截面 \(s(p)\) 与 \(s(q)\) 不能直接相减。普通偏导暗中使用了直积丛中“各点采用同一组基”的识别；一般向量丛没有这种天然识别。联络正是补上这一信息的微分规则。

\begin{definitionplain}[联络]
联络 \(\nabla\) 把向量场 \(X\) 和截面 \(s\) 送到截面 \(\nabla_Xs\)，并满足
\[
\nabla_{fX+gY}s=f\nabla_Xs+g\nabla_Ys,
\]
\[
\nabla_X(as+bt)=a\nabla_Xs+b\nabla_Xt,
\qquad
\nabla_X(fs)=X[f]s+f\nabla_Xs,
\]
其中 \(f,g\) 是光滑函数，\(a,b\in\mathbb C\)。第一式说明它对求导方向是函数线性的；最后一式是 Leibniz 法则。
\end{definitionplain}

在局部标架中写 \(s=e_a\psi^a\)。定义系数 \(1\)-形式 \(\Gamma^a{}_b\) 使
\[
\nabla e_b=e_a\Gamma^a{}_b.
\]
于是
\begin{equation}
\nabla s=e_a\bigl(\dd\psi^a+\Gamma^a{}_b\psi^b\bigr),
\qquad
\nabla\psi=\dd\psi+\Gamma\psi.
\label{eq:local-connection-primer}
\end{equation}
矩阵 \(\Gamma=\Gamma_\mu\dd x^\mu\) 不是额外的物理向量；它记录所选局部基在移动时怎样变化。

若更换局部标架 \(e'=eU\)，同一截面的坐标满足 \(\psi'=U^{-1}\psi\)。几何对象 \(\nabla s\) 不能因换基而改变，所以其坐标也应满足
\(\nabla'\psi'=U^{-1}\nabla\psi\)。将式 \eqnrefs{eq:local-connection-primer} 两边展开，得到
\[
\dd(U^{-1}\psi)+\Gamma'U^{-1}\psi
=U^{-1}(\dd\psi+\Gamma\psi).
\]
再用 \(\dd U^{-1}=-U^{-1}(\dd U)U^{-1}\)，比较 \(\psi\) 的系数，便有
\begin{equation}
\Gamma'=U^{-1}\Gamma U+U^{-1}\dd U.
\label{eq:connection-transform-primer}
\end{equation}
非齐次项 \(U^{-1}\dd U\) 说明联络系数本身不是张量。这不是缺陷：正是这个额外项抵消了 \(\dd\psi\) 换基时产生的导数项，使 \(\nabla\psi\) 按齐次方式变换。

沿曲线 \(\gamma:[0,1]\to M\) 的截面可写成 \(s(t)=e_a(\gamma(t))\psi^a(t)\)。称它被平行移动，是指
\[
\nabla_{\dot\gamma}s=0,
\qquad
\frac{\dd\psi}{\dd t}
+\Gamma_\mu(\gamma(t))\dot\gamma^\mu(t)\psi=0.
\]
给定 \(s(0)\) 后，这是线性常微分方程，因而局部存在唯一解。其解写成路径有序指数
\[
\psi(1)=
\mathcal P\exp\!\left[-\int_\gamma\Gamma\right]\psi(0).
\]
若 \(\Gamma\) 在不同点不对易，路径排序不可省略。闭路上的平行移动给出纤维到自身的线性变换，称为整体回转。

连续求两次协变导数时，二阶普通导数的对称部分在反对称化中抵消，剩下的量定义为曲率。对矩阵值形式，令协变外微分为
\[
D\alpha=\dd\alpha+\Gamma\wedge\alpha
\]
（若 \(\alpha\) 还带有右作用指标，相应再加右侧项）。作用在截面上可直接计算：
\[
D^2\psi
=\dd(\dd\psi+\Gamma\psi)
+\Gamma\wedge(\dd\psi+\Gamma\psi)
=(\dd\Gamma+\Gamma\wedge\Gamma)\psi.
\]
因此曲率二形式为
\begin{equation}
F=\dd\Gamma+\Gamma\wedge\Gamma,
\qquad
F_{\mu\nu}
=\partial_\mu\Gamma_\nu-
\partial_\nu\Gamma_\mu+[\Gamma_\mu,\Gamma_\nu].
\label{eq:curvature-primer}
\end{equation}
等价地，若坐标向量场对易，则
\[
[\nabla_\mu,\nabla_\nu]s=F_{\mu\nu}s.
\]
这给出了曲率的操作意义：沿两个无穷小方向平行移动，交换次序后的差异由 \(F_{\mu\nu}\) 控制。

把式 \eqnrefs{eq:connection-transform-primer} 代入式 \eqnrefs{eq:curvature-primer}，并使用 \(\dd^2=0\)，交叉项成对抵消，得到
\begin{equation}
F'=U^{-1}FU.
\label{eq:curvature-transform-primer}
\end{equation}
所以联络矩阵依赖标架，而曲率按共轭齐次变换。特别地，\(F=0\) 与否不依赖标架。继续计算协变外微分还得到 Bianchi 恒等式
\[
DF=\dd F+\Gamma\wedge F-F\wedge\Gamma=0.
\]

秩一情形中各矩阵退化为复数，\(\Gamma\wedge\Gamma=0\)，于是 \(F=\dd\Gamma\)。若局部规范变换写成 \(U=\ee^{\ii\chi}\)，则
\[
\Gamma'=\Gamma+\ii\dd\chi,
\qquad F'=F.
\]
这与电磁势改变而场强不变的结构相同。对秩大于一的丛，交换子项不能消失，曲率在换基下发生共轭而不是保持每个矩阵元不变。

分子问题中，参数 \(R\) 取代流形坐标，局部电子本征态 \(\{|u_a(R)\rangle\}\) 取代标架。矩阵元
\[
\Gamma_{\mu,ab}(R)=
\langle u_a(R)|\partial_\mu u_b(R)\rangle
\]
正是本征态丛上的联络。改变本征态基底会按式 \eqnrefs{eq:connection-transform-primer} 改变 \(\Gamma\)，但曲率、闭路整体回转以及由它们产生的可观测干涉不依赖这种选择。这样，后文的 Berry 相位、非 Abel 几何矩阵和 Born--Huang 协变核方程，都只是这里定义的同一套向量丛几何在不同物理系统中的具体实现。
